\documentclass[12pt,a4paper]{article}
\usepackage[utf8]{inputenc}
\usepackage[T1]{fontenc}
\usepackage{amsmath,amssymb}
\usepackage[margin=2.5cm]{geometry}
\usepackage{hyperref}

\hypersetup{
  pdfauthor={David Alfyorov, Igor Shnyukov},
  pdftitle={Erratum: Solar system and laboratory tests of spectral causal theory},
}

\title{\textbf{Erratum: Solar system and laboratory tests\\
of spectral causal theory}}
\author{David Alfyorov, Igor Shnyukov}
\date{April 2026}

\begin{document}
\maketitle

This erratum corrects two issues in Ref.~[1]
(doi:10.22541/au.177524450.03515205/v1).

\paragraph{1. Scalar Yukawa term in the modified Newtonian potential.}
Equation~(12) of Ref.~[1] gives the modified Newtonian potential as
\begin{equation}
  \frac{V(r)}{V_N(r)} = 1 - \frac{4}{3}\,e^{-m_2 r} + \frac{1}{3}\,e^{-m_0 r},
\end{equation}
with the scalar mass $m_0 = \Lambda\sqrt{6}$ at $\xi = 0$.
This formula uses the \emph{local} (Yukawa) approximation, in which
both the spin-2 and scalar poles of the propagator contribute.

Subsequent analysis~[2] has shown that the scalar propagator
denominator $\Pi_s(z,\xi) > 1$ for all positive Euclidean momenta
$z > 0$ and all $\xi$: the scalar pole does not exist in the full
nonlocal theory. The ``scalar mass'' $m_0 = \Lambda\sqrt{6}$ is an
artifact of the local approximation and does not correspond to a
physical propagator pole. The correct potential is the single-Yukawa
form
\begin{equation}
  \frac{V_{\mathrm{SCT}}(r)}{V_N(r)} = 1 - \frac{4}{3}\,e^{-m_2 r},
\end{equation}
where $m_2 = 1.554\,\Lambda$ (from the exact zero of $\Pi_{\mathrm{TT}}$,
correcting the local-approximation value $m_2^{\mathrm{loc}} = 2.148\,\Lambda$).

\paragraph{2. Torsion-balance bound.}
The bound $\Lambda > 2.38$~meV in Table~2 of Ref.~[1] was derived
using the double-Yukawa potential and the local spin-2 mass.
Using the corrected single-Yukawa formula with the exact mass
$m_2 = 1.554\,\Lambda$ and the Lee~et~al.\ (2020) exclusion curve
at $|\alpha| = 4/3$, the dedicated bound becomes
$\Lambda > 3.53$~meV.

The PPN parameters $\gamma - 1$ and $\beta - 1$ remain zero at
leading order (exponentially small corrections for $r \gg 1/\Lambda$),
so the Solar System conclusions of Ref.~[1] are unaffected.

\bigskip
\noindent [1] D.~Alfyorov and I.~Shnyukov,
``Solar system and laboratory tests of spectral causal theory,''
Zenodo (2025), doi:10.22541/au.177524450.03515205/v1.

\noindent [2] D.~Alfyorov and I.~Shnyukov,
``Absence of a real scalar graviton pole in spectral causal theory,''
in preparation (2026).

\end{document}
